\section{总结与延伸}

\subsection{整场讲座总结表}
\begin{table}[H]
\centering
\begin{tabular}{p{2.8cm}p{4.1cm}p{4.6cm}}
\toprule
\textbf{主题} & \textbf{讲者核心观点} & \textbf{对 Agent 研究的含义} \\
\midrule
开放性趋势 & 能力提升与访问下降同时发生 & 研究边界被访问权限直接塑形 \\
三层访问框架 & API / Open Weights / Open Source 分层清晰 & 不同层级对应不同可验证实验能力 \\
API Agent 工程 & LLM 充当执行控制器 & 重点转向记忆、反思、轨迹治理 \\
Benchmark 与实证 & 可评测体系是进步前提 & 从 demo 叙事转向可复现比较 \\
开放定义与治理 & 术语边界影响政策与协作 & “可下载”不等于“可复现” \\
长期路线 & 技术能力与基础设施共进化 & 开放科学需要标准、工具链与算力协同 \\
\bottomrule
\end{tabular}
\caption{Lecture 03 的核心论点压缩}
\end{table}

\subsection{延伸阅读}
\begin{itemize}
    \item Bommasani et al., \textit{On the Opportunities and Risks of Foundation Models}, 2021.
    \item Liang et al., \textit{Holistic Evaluation of Language Models (HELM)}, 2022.
    \item Agent evaluation and reproducibility 相关基准论文（含 MLAgentBench 系列工作）。
    \item OSI 关于 AI 模型开放定义的持续讨论与草案更新。
    \item 多 Agent 工程框架实践：AutoGen / LangGraph / CrewAI 的设计比较。
\end{itemize}

\subsection{收束观点}
\textit{``Access shapes research.''} 这句看似简单的话，实际上定义了 Agent 时代研究者的现实边界。  
如果希望在 Foundation Model 与 Agent 方向形成可持续科学进展，社区必须同时推进模型能力、开放访问与可复现基础设施，而不是只追逐单点 SOTA。

%% --- Fin del contenido --- %%

\end{document}
